% Unordered list
\begin{itemize}
    \item First item
    \item Second item
    \item Third item
    \item Fourth item
    \item Fifth item
\end{itemize}
% To change label use \item[label]

% Ordered list
\begin{enumerate}
    \item First item
    \item Second item
    \item Third item
    \item Fourth item
    \item Fifth item
\end{enumerate}

% Nested list
\begin{enumerate}
   \item First item
   \begin{itemize}
        \item First sub-item
        \item Second sub-item
        \begin{description}
        \item[Note:] Describing something
        \end{description}
   \end{itemize}
   \item Second item
\end{enumerate}

% Nested enumerated list (also possible with itemize)
\begin{enumerate}
    \item First level item
    \item First level item
    \begin{enumerate}
        \item Second level item
        \item Second level item
        \begin{enumerate}
            \item Third level item
            \item Third level item
            \begin{enumerate}
                \item Fourth level item
                \item Fourth level item
            \end{enumerate}
        \end{enumerate}
    \end{enumerate}
 \end{enumerate}

% Description
\begin{description}
   \item This is an entry \textit{without} a label.
   \item[Something short] A short one-line description.
   \item[Something long] A much longer description. Lorem ipsum dolor sit amet,
                         consetetur sadipscing elitr, sed diam nonumy eirmod
                         tempor invidunt ut labore et dolore magna aliquyam erat,
                         sed diam voluptua. At vero eos et accusam et justo duo
                         dolores et ea rebum. Stet clita kasd gubergren, no sea
                         takimata sanctus est Lorem ipsum dolor sit amet.
\end{description}